\documentclass[11pt,letterpaper]{article}

% Paquetes
\usepackage{fontspec}
\usepackage[utf8]{inputenc}
% Babel omitido intencionalmente: rompe \tableofcontents
% (su macro \babel@aux redefine \@writefile y vacía el TOC).
% xelatex + fontspec maneja UTF-8 nativo, así que tildes/ñ siguen funcionando.
\usepackage{graphicx}
\usepackage{xcolor}
\usepackage{geometry}
\usepackage[most]{tcolorbox}
\usepackage{booktabs}
\usepackage{longtable}
\usepackage{array}
\usepackage{enumitem}
\usepackage{fancyhdr}
\usepackage{parskip}
\usepackage[hidelinks]{hyperref}
\usepackage{float}
\usepackage{titlesec}
\usepackage{amssymb}
\usepackage{textcomp}
\usepackage{eso-pic}        % full-page background (cream en toda la hoja)

\geometry{
    letterpaper,
    left=1.8cm, right=1.8cm,
    top=2.2cm, bottom=2.2cm,
    headheight=14pt,
}

% Colores Consultorio Las Gaviotas — Concierge médico premium
% Paleta: navy / cream / gold / coral / ink
\definecolor{navy}{HTML}{0F1D3D}
\definecolor{navy-600}{HTML}{1A2C54}
\definecolor{navy-800}{HTML}{08122A}
\definecolor{navy-50}{HTML}{E8EBF3}
\definecolor{cream}{HTML}{F5F1E8}
\definecolor{surface}{HTML}{FAF7F0}
\definecolor{gold}{HTML}{A07C3E}
\definecolor{gold-600}{HTML}{8A6732}
\definecolor{gold-100}{HTML}{F1EAD9}
\definecolor{gold-50}{HTML}{F7F1E1}
\definecolor{coral}{HTML}{C8624A}
\definecolor{coral-600}{HTML}{A84A35}
\definecolor{coral-100}{HTML}{F7E1D8}
\definecolor{coral-50}{HTML}{FCEFEB}
\definecolor{ink}{HTML}{1A1A1A}
\definecolor{ink-2}{HTML}{5A5A5A}
\definecolor{muted}{HTML}{5A5A5A}
\definecolor{line}{HTML}{E6DFD1}
\definecolor{line-soft}{HTML}{F6F1E4}
\definecolor{paper}{HTML}{F5F1E8}
\definecolor{success}{HTML}{1F6E4E}
\definecolor{success-soft}{HTML}{D8E8E0}
\definecolor{warning}{HTML}{B47820}
\definecolor{warning-soft}{HTML}{F4EAD2}
\definecolor{danger}{HTML}{A4324A}
\definecolor{danger-soft}{HTML}{F0D8DF}
% Aliases para compatibilidad con código previo
\definecolor{brand}{HTML}{0F1D3D}
\definecolor{brand-50}{HTML}{E8EBF3}
\definecolor{brand-100}{HTML}{DDE4F0}
\definecolor{brand-700}{HTML}{1A2C54}
\definecolor{brand-800}{HTML}{08122A}
\definecolor{accent}{HTML}{A07C3E}

% Estilo del cuerpo
\setlength{\parindent}{0pt}
\setlength{\parskip}{0.6em}
\renewcommand{\baselinestretch}{1.18}

% Fuentes — Concierge médico premium
% Display: Noto Serif Display (peso semibold para títulos)
% Body: Noto Sans (peso regular)
% Mono: DejaVu Sans Mono (con box-drawing para árboles ASCII)
\setmainfont{Noto Serif Display}[Scale=1.0]
\setsansfont{Noto Sans}[Scale=1.0]
\setmonofont{DejaVu Sans Mono}[Scale=0.9]

% Page color: cream (full-bleed, cubre TODA la hoja, no solo el text area)
\AddToShipoutPictureBG*{\AtPageLowerLeft{\color{cream}\rule{\paperwidth}{\paperheight}}}

% Hyperlinks
\hypersetup{
    colorlinks=true,
    linkcolor=navy,
    urlcolor=gold,
    citecolor=navy,
    pdftitle={Sistema de Información Automatizado para la Gestión de Registro y Orden de Pacientes — Consultorio Las Gaviotas},
    pdfauthor={Consultorio Las Gaviotas},
    pdfsubject={Sistema de Información Automatizado para la Gestión de Registro y Orden de Pacientes en C.A Consultorio Médico Las Gaviotas, de Barcelona Estado Anzoátegui},
    pdfkeywords={Consultorio Las Gaviotas, RUP, UML, PostgreSQL, Bun, Astro, Concierge Médico},
}

% Cabeceras y pie — branding brass-plate
\pagestyle{fancy}
\fancyhf{}
\fancyhead[L]{\footnotesize\sffamily\bfseries\color{gold} Consultorio Las Gaviotas}
\fancyhead[R]{\footnotesize\color{ink-2} \thetitle}
\fancyfoot[L]{\footnotesize\sffamily\itshape\color{ink-2} Sistema de Información Automatizado · Barcelona, Anzoátegui}
\fancyfoot[C]{\footnotesize\sffamily\color{ink-2} \thepage\ /\ \pageref{LastPage}}
\fancypagestyle{plain}{%
  \fancyhf{}%
  \fancyhead[L]{\footnotesize\sffamily\bfseries\color{gold} Consultorio Las Gaviotas}%
  \fancyhead[R]{\footnotesize\color{ink-2} \thetitle}%
  \fancyfoot[L]{\footnotesize\sffamily\itshape\color{ink-2} Sistema de Información Automatizado · Barcelona, Anzoátegui}%
  \fancyfoot[C]{\footnotesize\sffamily\color{ink-2} \thepage\ /\ \pageref{LastPage}}%
}
\renewcommand{\headrulewidth}{0pt}
\renewcommand{\footrulewidth}{0pt}

% Estilo de títulos — Spectral/Noto Serif Display en navy con rule gold
\titleformat{\section}
  {\sffamily\bfseries\Huge\color{navy}}
  {}{0pt}{}[\vspace{-0.2em}\textcolor{gold}{\rule{\textwidth}{1.5pt}}\vspace{-0.2em}]
\titleformat{\subsection}
  {\sffamily\bfseries\LARGE\color{navy-600}}
  {}{0pt}{}
\titleformat{\subsubsection}
  {\sffamily\bfseries\Large\color{navy}}
  {}{0pt}{}
\titleformat{\paragraph}[runin]
  {\sffamily\bfseries\color{ink-2}}
  {}{0pt}{}

\titlespacing*{\section}{0pt}{1.5em}{0.6em}
\titlespacing*{\subsection}{0pt}{1.2em}{0.4em}
\titlespacing*{\subsubsection}{0pt}{1em}{0.3em}

% Cajas de colores concierge
\newtcolorbox{notebox}[1][Nota]{
  colback=gold-50, colframe=gold, fonttitle=\sffamily\bfseries,
  title=#1, breakable, enhanced
}
\newtcolorbox{warnbox}[1][Atención]{
  colback=warning-soft, colframe=warning, fonttitle=\sffamily\bfseries,
  title=#1, breakable, enhanced
}
\newtcolorbox{dangerbox}[1][Error]{
  colback=danger-soft, colframe=danger, fonttitle=\sffamily\bfseries,
  title=#1, breakable, enhanced
}
\newtcolorbox{successbox}[1][Listo]{
  colback=success-soft, colframe=success, fonttitle=\sffamily\bfseries,
  title=#1, breakable, enhanced
}

% Code blocks
\usepackage{tabularx}
\usepackage{adjustbox}
\usepackage{listings}
\usepackage{fancyvrb}
\lstset{
  basicstyle=\ttfamily\small,
  backgroundcolor=\color{surface},
  frame=single,
  framerule=0pt,
  framesep=8pt,
  xleftmargin=10pt, xrightmargin=10pt,
  breaklines=true,
  breakatwhitespace=true,
  numbers=left,
  numberstyle=\tiny\color{ink-2},
  numbersep=8pt,
  keywordstyle=\color{navy}\bfseries,
  commentstyle=\color{ink-2}\itshape,
  stringstyle=\color{coral},
  showstringspaces=false,
  tabsize=2,
  literate={á}{{\'a}}1 {é}{{\'e}}1 {í}{{\'i}}1 {ó}{{\'o}}1 {ú}{{\'u}}1 {ñ}{{\~n}}1,
}

% Last page label
\usepackage{lastpage}

\begin{document}
\begin{titlepage}
\thispagestyle{empty}
\centering
\vspace*{3cm}
\color{navy}\rule{0.7\textwidth}{1.2pt}\\[1.5em]
{\fontsize{36}{42}\selectfont\sffamily\bfseries\color{navy} Consultorio}\\[0.2em]
{\fontsize{42}{48}\selectfont\itshape\color{gold} Las Gaviotas}\\[1em]
\color{gold}\rule{0.7\textwidth}{1.2pt}\\[1.5em]
{\fontsize{22}{26}\selectfont\sffamily\color{navy-600} Índice de Entrega}\\[0.6em]
{\fontsize{14}{18}\selectfont\itshape\color{ink-2} Sistema de Información Automatizado para la Gestión de Registro}\\[0.2em]
{\fontsize{14}{18}\selectfont\itshape\color{ink-2} y Orden de Pacientes — Barcelona, Estado Anzoátegui}\\[2em]
{\fontsize{11}{14}\selectfont\sffamily\color{gold} \MakeUppercase{Análisis y Diseño de Sistemas · Framework RUP}}\\[0.3em]
{\fontsize{11}{14}\selectfont\sffamily\color{ink-2} \MakeUppercase{Concierge Médico Premium}}\\[3em]
\color{navy}\rule{0.5\textwidth}{0.6pt}\\[1em]
{\fontsize{12}{14}\selectfont\sffamily\color{ink-2} Fecha: 29 de julio de 2026}
\end{titlepage}
\thispagestyle{plain}
\renewcommand{\contentsname}{Índice}
\tableofcontents
\newpage
\section{Consultorio Las Gaviotas --- Índice de Entrega Final}

\textbf{Materia:} Análisis y Diseño de Sistemas
\textbf{Framework:} RUP (Rational Unified Process)
\textbf{Stack:} Bun + Elysia + PostgreSQL 16 + Astro 5 + Tailwind CSS 4
\textbf{Tipo de entrega:} Sistema de gestión para el Consultorio Médico Las Gaviotas (Barcelona, Anzoátegui)

\hrulefill

\subsection{1. Estructura de la entrega}

\begin{itemize}[leftmargin=*]
  \item \textbf{INDICE_ENTREGA.pdf} --- este documento.
  \item \textbf{01-documentos/}
  \item \texttt{vision.pdf} --- Documento de Visión.
  \item \texttt{catalogos.pdf} --- Catálogos (requisitos, excepciones, normas).
  \item \texttt{especificaciones-tecnicas.pdf} --- Especificaciones Técnicas.
  \item \textbf{02-diagramas/}
  \item \texttt{prototipo-ui.pdf} --- Prototipos UI estáticos.
  \item \texttt{casos-uso.pdf} --- Modelo de Casos de Uso.
  \item \texttt{diagrama-clases.pdf} --- Diagrama de Clases.
  \item \texttt{modelo-fisico.pdf} --- ER / Modelo Físico.
  \item \texttt{dsi-avanzado.pdf} --- Arquitectura + Realización CU.
  \item \textbf{03-software/}
  \item \texttt{prototipo-arquitectonico.pdf} --- Prototipo Arquitectónico.
  \item \texttt{README\_INSTALACION.pdf} --- Manual de instalación.
  \item \texttt{consultorio-gaviotas-codigo/} --- Código fuente completo.
\end{itemize}

\hrulefill

\subsection{2. Cómo usar esta entrega}

\subsubsection{2.1 Leer los documentos}

Los PDFs están listos para imprimir o visualizar. Cada uno incluye:
\begin{itemize}[leftmargin=*]
  \item Portada con el título del documento y metadatos de la materia.
  \item Numeración de páginas en el footer.
  \item Títulos jerárquicos con estilo académico.
\end{itemize}

\subsubsection{2.2 Revisar el software}


\begin{lstlisting}
cd entrega/03-software/consultorio-gaviotas-codigo
docker compose -f deploy/docker-compose.yml up -d
# Esperar 30s a que termine el seed.
# Abrir http://localhost:4321
# Login: admin / admin123
\end{lstlisting}



\subsubsection{2.3 Revisar el prototipo arquitectónico}


\begin{lstlisting}
cd entrega/03-software/consultorio-gaviotas-codigo/apps/backend
bun run prototype/consulta-flow.ts
\end{lstlisting}



\hrulefill

\subsection{3. Mapeo con los puntos de la entrega}

\begin{table}[H]
\centering
\small
\renewcommand{\arraystretch}{1.3}
\begin{tabularx}{\linewidth}{>{\RaggedRight\arraybackslash}X>{\RaggedRight\arraybackslash}X}
\toprule
\textbf{\textcolor{brand-800}{Punto de la entrega}} & \textbf{\textcolor{brand-800}{Archivo(s)}} \\
\midrule
\textbf{1. Documento de Visión} & \texttt{01-documentos/vision.pdf} \\
\textbf{1. Catálogos} (requisitos, excepciones, normas) & \texttt{01-documentos/catalogos.pdf} \\
\textbf{1. Especificaciones Técnicas} (entorno, seguridad, plan de pruebas, instalación) & \texttt{01-documentos/especificaciones-tecnicas.pdf} \\
\textbf{2. Prototipo UI estático} & \texttt{02-diagramas/prototipo-ui.pdf} \\
\textbf{2. Modelo de Casos de Uso} (con detalle en UC-06) & \texttt{02-diagramas/casos-uso.pdf} + \texttt{02-diagramas/dsi-avanzado.pdf} \\
\textbf{2. Diagrama de Clases} (Paciente, Cita, Historia clínica) & \texttt{02-diagramas/diagrama-clases.pdf} \\
\textbf{2. Modelo Físico de Datos} & \texttt{02-diagramas/modelo-fisico.pdf} \\
\textbf{2. DSI Avanzado} (arquitectura + realización de CU) & \texttt{02-diagramas/dsi-avanzado.pdf} \\
\textbf{3. Prototipo Arquitectónico} & \texttt{03-software/prototipo-arquitectonico.pdf} + \texttt{apps/backend/src/prototype/consulta-flow.ts} \\
\textbf{3. Aplicación Final (Consultorio Las Gaviotas)} & \texttt{03-software/consultorio-gaviotas-codigo/} \\
\bottomrule
\end{tabularx}
\end{table}
\hrulefill

\subsection{4. Credenciales de acceso}

\begin{table}[H]
\centering
\small
\renewcommand{\arraystretch}{1.3}
\begin{tabularx}{\linewidth}{>{\RaggedRight\arraybackslash}X>{\RaggedRight\arraybackslash}X>{\RaggedRight\arraybackslash}X}
\toprule
\textbf{\textcolor{brand-800}{Usuario}} & \textbf{\textcolor{brand-800}{Contraseña}} & \textbf{\textcolor{brand-800}{Rol}} \\
\midrule
admin & admin123 & ADMIN \\
medico & medico123 & MEDICO \\
medico2 & medico123 & MEDICO \\
recep & recep123 & RECEPCION \\
recep2 & recep123 & RECEPCION \\
\bottomrule
\end{tabularx}
\end{table}
\begin{quote}\itshape Cambiar todas las contraseñas antes de producción.\end{quote}

\hrulefill

\subsection{5. Estado del proyecto}

Sistema funcional, 3 fases RUP completas (Inicio, Elaboración, Construcción). Desplegado en Railway con 1 servicio y 3 procesos (web Astro SSR, API Elysia, PostgreSQL).

\end{document}
